% Part: first-order-logic
% Chapter: syntax-and-semantics
% Section: terms-formulas

\documentclass[../../../include/open-logic-section]{subfiles}

\begin{document}

\olfileid{fol}{syn}{frm}

\olsection{ٹرم تے \printtoken{P}{formula}}

جدوں پہلے درجے دی اک زبان~$\Lang L$ دتی جاوے، تاں اسیں
زبان~$\Lang L$ دی بنیادی علامتاں توں بندیاں لکھتاں دی تعریف کر سکدے آں۔
ایہناں وچ خاص طور تے \emph{ٹرم} تے \emph{!!{formula}s} شامل نیں۔

\begin{defn}[ٹرم]
\ollabel{defn:terms}
\emph{ٹرماں} دا سیٹ~$\Trm[L]$، $\Lang L$ لئی، اینج
استقرائی طور تے تعریف کیتا جاندا اے:
\begin{enumerate}
\item ہر !!{variable} اک ٹرم اے۔
\item $\Lang L$ دا ہر !!{constant} اک ٹرم اے۔
\item جے $f$ اک $n$-مقامی !!{function} اے تے $t_1$, \dots, $t_n$
  ٹرم نیں، تاں $\Atom{f}{t_1, \ldots, t_n}$ وی اک ٹرم اے۔
\tagitem{limitClause}{
  ہور کوئی شے ٹرم نہیں اے۔}{}
\end{enumerate}
جیہڑے ٹرم وچ کوئی !!{variable}s نہ ہون، اوہ \emph{بند ٹرم} اے۔
\end{defn}

\begin{explain}
ساڈی زبان تے ٹرماں دی وضاحت وچ !!{constant}s علامتاں دی
وکھری قسم دے طور تے آؤندے نیں، پر اسیں انہاں نوں
صفر-مقامی !!{function}s وچ وی شامل کر سکدے ساں۔ فیر
ٹرماں دی تعریف دی دوجی شق دی لوڑ نہ رہندی۔ بس
$\Atom{f}{t_1, \ldots, t_n}$ نوں اکیلا $f$
سمجھنا پیندا اے، جے $n =
0$ ہووے۔
\end{explain}

\begin{defn}[فارمولے]
\ollabel{defn:formulas}
\emph{!!{formula}s} دا سیٹ~$\Frm[L]$، زبان~$\Lang L$ لئی،
اینج استقرائی طور تے تعریف کیتا جاندا اے:
\begin{enumerate}
\tagitem{prvFalse}{$\lfalse$ اک ایٹمی !!{formula} اے۔}{}

\tagitem{prvTrue}{$\ltrue$ اک ایٹمی !!{formula} اے۔}{}

\item جے $R$ اک $n$-مقامی !!{predicate} اے جو~$\Lang L$ دا اے، تے $t_1$, \dots,
  $t_n$، $\Lang L$ دے ٹرم نیں، تاں $\Atom{R}{t_1,\ldots, t_n}$ اک
  ایٹمی !!{formula} اے۔

\item جے $t_1$ تے $t_2$، $\Lang L$ دے ٹرم نیں، تاں $\Atom{\eq}{t_1, t_2}$
  اک ایٹمی !!{formula} اے۔

\tagitem{prvNot}{جے $!A$، !!a{formula} اے، تاں $\lnot !A$ وی
  !!a{formula} اے۔}{}

\tagitem{prvAnd}{جے $!A$ تے $!B$، !!{formula}s نیں، تاں $(!A \land
  !B)$ وی !!a{formula} اے۔}{}

\tagitem{prvOr}{جے $!A$ تے $!B$، !!{formula}s نیں، تاں $(!A \lor !B)$
  وی !!a{formula} اے۔}{}

\tagitem{prvIf}{جے $!A$ تے $!B$، !!{formula}s نیں، تاں $(!A \lif !B)$
  وی !!a{formula} اے۔}{}

\tagitem{prvIff}{جے $!A$ تے $!B$، !!{formula}s نیں، تاں $(!A \liff !B)$
  وی !!a{formula} اے۔}{}

\tagitem{prvAll}{جے $!A$، !!a{formula} اے تے $x$، !!a{variable} اے،
  تاں $\lforall[x][!A]$ وی !!a{formula} اے۔}{}

\tagitem{prvEx}{جے $!A$، !!a{formula} اے تے $x$، !!a{variable} اے،
  تاں $\lexists[x][!A]$ وی !!a{formula} اے۔}{}

\tagitem{limitClause}{ہور کوئی شے !!a{formula} نہیں اے۔}{}
\end{enumerate}
\end{defn}

\begin{explain}
ٹرماں دے سیٹ تے !!{formula}s دے سیٹ دیاں تعریفاں
\emph{استقرائی تعریفاں} نیں۔ بنیادی طور تے اسیں
!!{formula}s دا سیٹ لامتناہی مرحلیاں وچ بناؤندے آں۔
پہلے مرحلے وچ اسیں سارے ایٹمی فارمولیاں نوں فارمولے
قرار دیندے آں؛ ایہ تعریف دیاں پہلیاں کجھ صورتاں، یعنی
\iftag{prvTrue}{$\ltrue$, }{}%
\iftag{prvFalse}{$\lfalse$, }{}%
$\Atom{R}{t_1,\dots,t_n}$ تے $\Atom{\eq}{t_1,t_2}$
والیاں صورتاں دے مطابق اے۔ ایس لئی ”ایٹمی
!!{formula}“ دا مطلب ایس روپ دا کوئی وی !!{formula} اے۔

تعریف دیاں ہور صورتاں پہلے بنائے !!{formula}s توں
نویں !!{formula}s بناؤن دے قاعدے دیندیاں نیں۔
دوجے مرحلے وچ اسیں اوہناں نال ایٹمی !!{formula}s توں
!!{formula}s بنا سکدے آں۔ تیجے مرحلے وچ اسیں ایٹمی
فارمولیاں تے دوجے مرحلے وچ ملے فارمولیاں توں نویں
فارمولے بناؤندے آں، تے ایویں اگے۔ !!{formula} صرف
اوہ شے اے جیہڑی آخرکار کسے ایہو جہے مرحلے اُتے
بنائی جاوے؛ ہور کوئی شے نہیں۔
\end{explain}

رواجاً اسیں $\eq$ نوں اوہدیاں دو دلیلاں دے وچ لکھدے تے
قوساں چھڈ دیندے آں: $\eq[t_1][t_2]$،
$\Atom{\eq}{t_1,t_2}$ دا مخفف اے۔ اس توں علاوہ
$\lnot \Atom{\eq}{t_1,t_2}$ نوں $\eq/[t_1][t_2]$
لکھ کے مختصر کردے آں۔ جدوں فارمولا $(!B \ast !C)$،
جیہڑا $!B$ تے $!C$ توں دو-مقامی رابطے~$\ast$
نال بنیا اے، لکھدے آں تاں اکثر سب توں باہری قوساں دی جوڑی
چھڈ کے صرف~$!B \ast !C$ لکھدے آں۔

\begin{intro}
منطق دیاں کجھ کتاباں وچ !!{variable}~$x$ دا~$!A$
وچ آنا لازمی منیا جاندا اے، تاں جو
\iftag{prvEx}{$\lexists[x][!A]$ }{}%
\iftag{notprvEx,notprvAll}{}{تے }%
\iftag{prvAll}{$\lforall[x][!A]$ }{}%
نوں
\iftag{notprvEx,notprvAll}{!!a{formula}}{!!{formula}s}
منیا جائے۔ ایہ شرط نہ رکھن نال کوئی خرابی نہیں آؤندی،
تے کم آسان ہو جاندا اے۔
\end{intro}

\begin{tagblock}{defNot,defOr,defAnd,defIf,defIff,defTrue,defFalse,defEx,defAll}
\begin{defn}
تعریف شدہ عملگراں نال بنے فارمولیاں نوں اینج سمجھیا جاوے:

\begin{tagenumerate}{defTrue,defFalse,defNot,defOr,defAnd,defIf,defIff,defEx,defAll}

\tagitem{defTrue}{$\ltrue$،
  \iftag{prvFalse}{$\lnot\lfalse$}{$(!A \lor \lnot !A)$، کسے
    مقرر ایٹمی !!{formula}~$!A$ لئی} دا مخفف اے۔}{}

\tagitem{defFalse}{$\lfalse$،
  \iftag{prvTrue}{$\lnot\ltrue$}{$(!A \land \lnot !A)$، کسے
    مقرر ایٹمی !!{formula}~$!A$ لئی} دا مخفف اے۔}{}

\tagitem{defNot}{$\lnot !A$، $!A \lif \lfalse$ دا مخفف اے۔}{}

\tagitem{defOr}{$!A \lor !B$،
  \iftag{prvAnd}{$\lnot(\lnot !A \land \lnot !B)$}{$\lnot !A \lif
    !B$} دا مخفف اے۔}{}

\tagitem{defAnd}{$!A \land !B$،
  \iftag{prvOr}{$\lnot(\lnot !A \lor \lnot !B)$}{$\lnot (!A \lif
    \lnot !B)$} دا مخفف اے۔}{}

% PNB-SYN-001 ماخذ دی درستی: defIf دی پہلی متبادل عبارت وچ بند ہُندی قوس توں پہلاں کھلدی قوس ماخذ وچ رہ گئی؛ ایتھے پچھلے PNBPL-002 وانگوں جوڑی پوری کیتی اے۔
\tagitem{defIf}{$!A \lif !B$،
  \iftag{prvOr}{$(\lnot !A \lor !B)$}{$\lnot (!A \land \lnot !B)$} دا مخفف اے۔}{}

\tagitem{defIff}{$!A \liff !B$، $(!A \lif !B) \land (!B
  \lif !A)$ دا مخفف اے۔}{}

\tagitem{defAll}{$\lforall[x][!A]$، $\lnot\lexists[x][\lnot !A]$ دا مخفف اے۔}{}

\tagitem{defEx}{$\lexists[x][!A]$، $\lnot\lforall[x][\lnot !A]$ دا مخفف اے۔}{}
\end{tagenumerate}
\end{defn}
\end{tagblock}

جے اسیں کسے خاص ورتوں دی زبان وچ کم کریے، تاں دو-مقامی
!!{predicate}s تے !!{function}s نوں اکثر اوہناں دے متعلقہ
ٹرماں دے وچ لکھدے آں، مثلاً حساب دی زبان وچ $t_1 < t_2$
تے $(t_1 + t_2)$، تے سیٹ نظریے دی زبان وچ $t_1 \in t_2$۔
حساب دی زبان دا جانشین فنکشن رواجاً اپنی دلیل دے
\emph{بعد} وی لکھیا جاندا اے:~$t'$۔ پر باقاعدہ طور تے
ایہ صرف $\Atom{\Obj A^2_0}{t_1, t_2}$،
$\Obj f^2_0(t_1, t_2)$، $\Atom{\Obj A^2_0}{t_1, t_2}$
تے $f^1_0(t)$ دے رواجی مخفف نیں، بالترتیب۔

\begin{defn}[نحوی یکسانیت]
علامت $\ident$، علامتاں دیاں لڑیاں وچ نحوی یکسانیت
ظاہر کردی اے، یعنی $!A \ident !B$ اُدوں تے صرف
اُدوں جدوں $!A$ تے $!B$ اکّو لمبائی دیاں علامتی
لڑیاں ہون تے ہر تھاں اُتے اکّو علامت رکھدیاں ہون۔
\end{defn}

علامت $\ident$ دے دونیں پاسے جوڑ نال ملیاں لڑیاں
آ سکدیاں نیں؛ مثلاً $!A \ident (!B \lor !C)$ دا
مطلب اے کہ علامتاں دی لڑی~$!A$ اوہی لڑی اے
جیہڑی اک کھلدی قوس، لڑی $!B$، علامت $\lor$،
لڑی~$!C$ تے اک بند ہُندی قوس نوں ایس ترتیب نال
جوڑ کے ملدی اے۔ جے ایہ گل ہووے، تاں اسیں جاندے
آں کہ~$!A$ دی پہلی علامت کھلدی قوس اے، $!A$
وچ $!B$ اک ذیلی لڑی دے طور تے شامل اے (دوجی
علامت توں شروع ہو کے)، اس ذیلی لڑی دے بعد~$\lor$
آؤندا اے، وغیرہ۔

جیویں ٹرم تے !!{formula}s بنیادی عنصراں توں
استقرائی تعریفاں راہیں بندے نیں، اسیں انہاں بارے
گلاں ثابت کرن لئی ایہ استقرائی اصول ورت سکدے آں۔

\begin{lem}[\emph{ٹرماں اُتے استقرا دا اصول}]
    \ollabel{lem:trmind}
    فرض کرو $\Lang L$ پہلے درجے دی زبان اے۔
    جے کوئی خاصیت~$P$ اینج ہووے کہ
    %
    \begin{enumerate}
        \item اوہ ہر !!{variable}~$v$ لئی درست ہووے،
        %
        \item اوہ ہر !!{constant}~$a$ لئی درست ہووے جو~$\Lang L$ دا اے، تے
        %
        \item اوہ $f(t_1,\dotsc,t_n)$ لئی وی درست ہووے جدوں
            اوہ $t_1$,~\dots, $t_n$ لئی درست ہووے تے $f$
            اک $n$-مقامی !!{function} ہووے جو~$\Lang L$ دا اے
    \end{enumerate}
    (ایہ من کے کہ $t_1$,~\dots, $t_n$، $\Lang{L}$ دے ٹرم نیں)،
    تاں $P$، $\Trm[L]$ دے ہر ٹرم لئی درست اے۔
\end{lem}

\begin{prob}
    \olref[fol][syn][frm]{lem:trmind} ثابت کرو۔
\end{prob}

\begin{lem}[\emph{!!{formula}s اُتے استقرا دا اصول}]
    \ollabel{thm:frmind}
    فرض کرو $\Lang L$ پہلے درجے دی زبان اے۔
    جے کوئی خاصیت~$P$ سارے ایٹمی !!{formula}s لئی
    درست ہووے، تے اینج وی ہووے کہ
    %
    \begin{enumerate}
        \tagitem{prvNot}{اوہ $\lnot !A$ لئی وی درست ہووے
            جدوں اوہ~$!A$ لئی درست ہووے؛}{}
        \tagitem{prvAnd}{اوہ $(!A \land !B)$ لئی وی درست ہووے
            جدوں اوہ $!A$ تے~$!B$ لئی درست ہووے؛}{}
        \tagitem{prvOr}{اوہ $(!A \lor !B)$ لئی وی درست ہووے
            جدوں اوہ $!A$ تے~$!B$ لئی درست ہووے؛}{}
        \tagitem{prvIf}{اوہ $(!A \lif !B)$ لئی وی درست ہووے
            جدوں اوہ $!A$ تے~$!B$ لئی درست ہووے؛}{}
        \tagitem{prvIff}{اوہ $(!A \liff !B)$ لئی وی درست ہووے
            جدوں اوہ $!A$ تے~$!B$ لئی درست ہووے؛}{}
        \tagitem{prvEx}{اوہ $\lexists[x][!A]$ لئی وی درست ہووے
            جدوں اوہ~$!A$ لئی درست ہووے؛}{}
        \tagitem{prvAll}{اوہ $\lforall[x][!A]$ لئی وی درست ہووے
            جدوں اوہ~$!A$ لئی درست ہووے؛}{}
  \end{enumerate}
  (ایہ من کے کہ $!A$ تے $!B$، $\Lang{L}$ دے !!{formula}s نیں)،
  تاں $P$، $\Frm[L]$ دے سارے فارمولیاں لئی درست اے۔
\end{lem}

\end{document}
